% Part: normal-modal-logic
% Chapter: tableaux
% Section: soundness

\documentclass[../../../include/open-logic-section]{subfiles}

\begin{document}

\olfileid{nml}{tab}{sou}

\olsection{\Log{K} માટેની યથાર્થતા}

\begin{editorial}
  યથાર્થતાનું આ પ્રમાણ શાસ્ત્રીય વિધાનતર્કનું યથાર્થતા-પ્રમાણ ફરી
  વાપરે છે: એટલે તે બધું શરૂઆતથી સાબિત કરે છે. સ્વતંત્ર રીતે
  વાંચી શકાય એવું યથાર્થતા-પ્રમાણ જોઈએ તો આ યોગ્ય છે. સામાન્ય
  ટેબ્લો માટેની યથાર્થતા પહેલાં જોઈ હોય તો આ પુનરાવર્તન લાગશે.
  આગળ જતાં સ્વતંત્ર આવૃત્તિ અને અમોડલ કિસ્સા પર આધારિત આવૃત્તિ
  વચ્ચે પસંદગી શક્ય બનાવવાની યોજના છે.
\end{editorial}

\begin{explain}
  ઉપસર્ગવાળા !!{tableau}s યથાર્થ છે તે બતાવવા આપણે સાબિત કરવું પડે કે
  \[
  \sFmla{\True}{!B_1}[1], \dots, \sFmla{\True}{!B_n}[1], \sFmla{\False}{!A}[1]
  \]
  માટે બંધ !!{tableau} હોય, તો $!B_1, \dots, !B_n \Entails !A$.
  પ્રતિવિધાન સાબિત કરવું વધુ સરળ છે: કોઈ $\mModel{M}$ અને
  જગત~$w$ માટે બધા $i=1$, \dots,~$n$ પર
  $\mSat{M}{!B_i}[w]$ હોય, પરંતુ
  $\mSat/{M}{!A}[w]$ હોય, તો કોઈ !!{tableau} બંધ થઈ શકે નહિ.
  \sourcecorrection{OLMD-052}{મૂળમાં પ્રતિનિદર્શનમાં $!A$ પણ સત્ય
  લખાયું છે; પ્રતિવિધાન અને આગળની દલીલ માટે તે અસત્ય હોવું જરૂરી
  છે, તેથી અસંતોષનું ચિહ્ન પુનઃસ્થાપિત કર્યું છે.}
  આવું પ્રતિનિદર્શન બતાવે છે કે !!{tableau}ની પ્રારંભિક ધારણાઓ
  સંતોષ્ય છે. પ્રમાણની રીત એ બતાવવાની છે કે !!a{tableau}ની કોઈ
  શાખા પરનાં બધાં ઉપસર્ગવાળાં !!{formula}s સંતોષ્ય હોય ત્યારે
  કોઈપણ નિયમ લાગુ કરવાથી ઓછામાં ઓછી એક વિસ્તૃત શાખા પણ
  સંતોષ્ય રહે છે. બંધ શાખાઓ અસંતોષ્ય હોવાથી, ઉપસર્ગવાળાં
  !!{formula}sના સંતોષ્ય ગણ માટેના દરેક !!{tableau}માં
  ઓછામાં ઓછી એક ખુલ્લી શાખા હશે.

  આ રીત મોડલ કિસ્સામાં લાગુ કરવા, ``સંતોષ્ય''ની વ્યાખ્યા
  મોડલ નિદર્શનો અને ઉપસર્ગો સુધી વિસ્તૃત કરવી પડે. તે થઈ જાય
  પછી પ્રમાણ સીધું છે.
\end{explain}

\begin{defn}
  $P$ ઉપસર્ગોનો ગણ હોય, એટલે $P \subseteq (\PosInt)^*
  \setminus \{\emptyseq\}$, અને $\mModel{M}$ નિદર્શન હોય.
  વિધેય~$f\colon P \to W$ એ~$\mModel{M}$માં~$P$નું
  \emph{અર્થઘટન} ત્યારે કહેવાય જ્યારે $\sigma$ અને $\sigma.n$
  બંને~$P$માં હોય ત્યારે $Rf(\sigma)f(\sigma.n)$ હોય.

  ઉપસર્ગો~$P$ના અર્થઘટનને અનુલક્ષીને વ્યાખ્યાયિત કરીએ છીએ:
  \begin{enumerate}
  \item $\mModel{M}$ એ $\sFmla{\True}{!A}[\sigma]$ને સંતોષે છે
    જો અને માત્ર જો
    $\mSat{M}{!A}[f(\sigma)]$.
  \item $\mModel{M}$ એ $\sFmla{\False}{!A}[\sigma]$ને સંતોષે છે
    જો અને માત્ર જો
    $\mSat/{M}{!A}[f(\sigma)]$.
  \end{enumerate}
\end{defn}

\begin{defn}
  $\Gamma$ ઉપસર્ગવાળાં !!{formula}sનો ગણ અને $P(\Gamma)$ તેમાં
  આવતા ઉપસર્ગોનો ગણ હોય. $f$ એ~$\mModel{M}$માં
  $P(\Gamma)$નું અર્થઘટન હોય, તો $\mModel{M}$ એ~$f$ને અનુલક્ષીને
  $\Gamma$ને સંતોષે છે, $\mSat{M}{\Gamma}[f]$, એમ ત્યારે કહીએ
  જ્યારે $\mModel{M}$ એ~$\Gamma$ના દરેક ઉપસર્ગવાળા
  !!{formula}ને~$f$ને અનુલક્ષીને સંતોષે. $\Gamma$
  \emph{સંતોષ્ય} છે જો અને માત્ર જો એવું નિદર્શન~$\mModel{M}$
  અને $P(\Gamma)$નું અર્થઘટન~$f$ હોય કે
  $\mSat{M}{\Gamma}[f]$.
\end{defn}

\begin{prop}
  જો $\Gamma$માં $\sFmla{\True}{!A}[\sigma]$ અને
  $\sFmla{\False}{!A}[\sigma]$ બંને હોય, કોઈ !!{formula}~$!A$
  અને ઉપસર્ગ~$\sigma$ માટે, તો $\Gamma$ અસંતોષ્ય છે.
\end{prop}

\begin{proof}
  એવું નિદર્શન~$\mModel{M}$ અને $P(\Gamma)$નું અર્થઘટન~$f$
  હોઈ શકે નહિ કે $\mSat{M}{!A}[f(\sigma)]$ અને
  $\mSat/{M}{!A}[f(\sigma)]$ બંને હોય.
\end{proof}

\begin{thm}[યથાર્થતા]
  \ollabel{thm:tableau-soundness}
  જો $\Gamma$ માટે બંધ !!{tableau} હોય, તો $\Gamma$ અસંતોષ્ય છે.
\end{thm}

\begin{proof}
!!a{tableau}ની શાખા પરનાં !!{signed formula}sનો ગણ સંતોષ્ય હોય
તો અને માત્ર તો શાખાને સંતોષ્ય કહીએ; અને ઓછામાં ઓછી એક સંતોષ્ય
શાખા ધરાવતું !!a{tableau} સંતોષ્ય કહેવાય.

નીચેનું બતાવીએ છીએ: સંતોષ્ય !!{tableau}ને અનુમાનના કોઈ એક નિયમથી
વિસ્તારતાં હંમેશા સંતોષ્ય !!{tableau} મળે છે. તેનાથી પ્રમેય
સાબિત થશે: કોઈપણ બંધ !!{tableau}, માત્ર~$\Gamma$માંથી લીધેલી
ધારણાઓવાળા !!{tableau}ને અનુમાનના નિયમો લાગુ કરવાથી મળે છે.
તેથી $\Gamma$ સંતોષ્ય હોય, તો તેના માટેનું દરેક !!{tableau}
સંતોષ્ય હશે. પરંતુ બંધ !!{tableau} દેખીતી રીતે સંતોષ્ય નથી:
તેની બધી શાખાઓ બંધ છે અને બંધ શાખાઓ અસંતોષ્ય છે.

ધારો કે આપણી પાસે સંતોષ્ય !!{tableau} છે, એટલે ઓછામાં ઓછી એક
સંતોષ્ય શાખાવાળું !!a{tableau}. અનુમાનનો નિયમ લાગુ કરવાથી
શાખામાં !!{signed formula}s ઉમેરાય છે અથવા શાખા બે ભાગમાં
વહેંચાય છે. જો !!{tableau}માં એવી સંતોષ્ય શાખા હોય જેને સંબંધિત
નિયમથી વિસ્તારી ન હોય, તો વિસ્તૃત !!{tableau}માં પણ તે
સંતોષ્ય જ રહે છે; તેથી વિસ્તૃત ટેબ્લો સંતોષ્ય છે. આમ માત્ર એવો
કિસ્સો ધ્યાનમાં લેવો પડે જેમાં નિયમ સંતોષ્ય શાખાને લાગુ થાય.

તે શાખા પરનાં !!{signed formula}sનો ગણ $\Gamma$ લો, અને નિયમ
જેને લાગુ પડે તે !!{signed formula}
$\sFmla{S}{!A}[\sigma] \in \Gamma$ લો. જો નિયમથી શાખા
વહેંચાતી ન હોય, તો વિસ્તૃત શાખા, એટલે નિયમના નિષ્કર્ષો સાથેનો
$\Gamma$, હજી સંતોષ્ય છે તે બતાવવું પડે. જો શાખા વહેંચાતી
હોય, તો મળતી બે શાખામાંથી ઓછામાં ઓછી એક સંતોષ્ય છે તે બતાવવું
પડે.
\sourcecorrection{OLMD-053}{મૂળમાં અહીં
\texttt{\string\tagfalse\{prvDiamond\}} લખીને ડાયમંડ-શાખાઓ
બંધ કરી દેવાય છે, જ્યારે આગળની સાબિતીમાં એ જ શાખાઓ જરૂરી છે.
આ અણધાર્યો ટૅગ-ફેરફાર દૂર કર્યો છે.}
પહેલાં શાખા ન વહેંચતા સંભવિત અનુમાનો ધ્યાનમાં લઈએ.
\begin{enumerate}
\item $\sFmla{\True}{\lnot !B}[\sigma] \in \Gamma$ને
  $\TRule{\True}{\lnot}$ લાગુ કરીને શાખા વિસ્તારીએ. ત્યારે
  વિસ્તૃત શાખામાં !!{signed formula}sનો ગણ $\Gamma \cup
  \{\sFmla{\False}{!B}[\sigma]\}$ છે. ધારો કે
  $\mSat{M}{\Gamma}[f]$. ખાસ કરીને,
  $\mSat{M}{\lnot !B}[f(\sigma)]$. તેથી
  $\mSat/{M}{!B}[f(\sigma)]$, એટલે~$\mModel{M}$ એ
  $\sFmla{\False}{!B}[\sigma]$ને~$f$ને અનુલક્ષીને સંતોષે છે.
\item $\sFmla{\False}{\lnot !B}[\sigma] \in \Gamma$ને
  $\TRule{\False}{\lnot}$ લાગુ કરીને શાખા વિસ્તારીએ:
  અભ્યાસપ્રશ્ન.
\item $\sFmla{\True}{!B \land !C}[\sigma] \in \Gamma$ને
  $\TRule{\True}{\land}$ લાગુ કરીને શાખા વિસ્તારીએ. તેનાથી
  શાખા પર બે નવાં !!{signed formula}s મળે છે:
  $\sFmla{\True}{!B}[\sigma]$ અને $\sFmla{\True}{!C}[\sigma]$.
  ધારો કે $\mSat{M}{\Gamma}[f]$, ખાસ કરીને
  $\mSat{M}{!B \land !C}[f(\sigma)]$. ત્યારે
  $\mSat{M}{!B}[f(\sigma)]$ અને $\mSat{M}{!C}[f(\sigma)]$.
  એટલે $\mModel{M}$ એ $\sFmla{\True}{!B}[\sigma]$ અને
  $\sFmla{\True}{!C}[\sigma]$ બંનેને~$f$ને અનુલક્ષીને સંતોષે છે.
\item $\sFmla{\False}{!B \lor !C}[\sigma] \in \Gamma$ને
  $\TRule{\False}{\lor}$ લાગુ કરીને શાખા વિસ્તારીએ:
  અભ્યાસપ્રશ્ન.\sourcecorrection{OLMD-054}{મૂળમાં આ ચિહ્નિત
  સૂત્રનો ઉપસર્ગ છૂટી ગયો છે; શાખાના અન્ય નિયમો સાથે સુસંગત
  રહે તે માટે $[\sigma]$ ઉમેર્યું છે.}
\item $\sFmla{\False}{!B \lif !C}[\sigma] \in \Gamma$ને
  $\TRule{\False}{\lif}$ લાગુ કરીને શાખા વિસ્તારીએ. તેનાથી
  શાખા પર બે નવાં !!{signed formula}s મળે છે:
  $\sFmla{\True}{!B}[\sigma]$ અને
  $\sFmla{\False}{!C}[\sigma]$. ધારો કે
  $\mSat{M}{\Gamma}[f]$, ખાસ કરીને
  $\mSat/{M}{!B \lif !C}[f(\sigma)]$. ત્યારે
  $\mSat{M}{!B}[f(\sigma)]$ અને
  $\mSat/{M}{!C}[f(\sigma)]$. એટલે
  $\mModel{M}, f$ એ $\sFmla{\True}{!B}[\sigma]$ અને
  $\sFmla{\False}{!C}[\sigma]$ બંનેને સંતોષે છે.

\iftag{prvBox}{%
\item $\sFmla{\True}{\Box !B}[\sigma] \in \Gamma$ને
  $\TRule{\True}{\Box}$ લાગુ કરીને શાખા વિસ્તારીએ:
  \iftag{probBox}{અભ્યાસપ્રશ્ન.}{તેનાથી કોઈ
    $\sigma.n \in P(\Gamma)$ માટે નવું !!{signed
      formula}~$\sFmla{\True}{!B}[\sigma.n]$ શાખા પર મળે છે
    (કારણ કે $\sigma.n$ વપરાયેલો હોવો જોઈએ). ધારો કે
    $\mSat{M}{\Gamma}[f]$, ખાસ કરીને
    $\mSat{M}{\Box !B}[f(\sigma)]$. $f$ ઉપસર્ગોનું અર્થઘટન છે
    અને $\sigma$, $\sigma.n \in P(\Gamma)$ બંને હોવાથી
    $Rf(\sigma)f(\sigma.n)$ છે. તેથી
    $\mSat{M}{!B}[f(\sigma.n)]$; એટલે $\mModel{M}, f$ એ
    $\sFmla{\True}{!B}[\sigma.n]$ને સંતોષે છે.}

\item $\sFmla{\False}{\Box !B}[\sigma] \in \Gamma$ને
  $\TRule{\False}{\Box}$ લાગુ કરીને શાખા વિસ્તારીએ:
  \iftag{probBox}{અભ્યાસપ્રશ્ન.}{તેનાથી નવું !!{signed
      formula}~$\sFmla{\False}{!B}[\sigma.n]$ મળે છે, જ્યાં
    $\sigma.n$ શાખા પર નવો ઉપસર્ગ છે, એટલે
    $\sigma.n \notin P(\Gamma)$.
    \sourcecorrection{OLMD-055}{મૂળમાં આ નિયમનું નવું સૂત્ર $!A$
    સાથે લખાયું છે; પૂર્વપક્ષ $\Box !B$ અને આગળની સાબિતી મુજબ
    અહીં $!B$ હોવું જોઈએ.}
    $\Gamma$ સંતોષ્ય હોવાથી એવી $\Struct{M}$ અને
    $P(\Gamma)$નું અર્થઘટન~$f$ છે કે
    $\mSat{M}{\Gamma}[f]$, ખાસ કરીને
    $\mSat/{M}{\Box !B}[f(\sigma)]$.
    $\Gamma \cup \{\sFmla{\False}{!B}[\sigma.n]\}$ સંતોષ્ય છે
    તે બતાવવાનું છે. તે માટે $P(\Gamma)
    \cup \{\sigma.n\}$નું અર્થઘટન નીચે મુજબ વ્યાખ્યાયિત કરીએ:

  $\mSat/{M}{\Box !B}[f(\sigma)]$ હોવાથી એવું $w \in W$ છે કે
  $Rf(\sigma)w$ અને $\mSat/{M}{!B}[w]$. $f'$ એ $f$ જેવું જ
  રાખો, માત્ર $f'(\sigma.n) = w$. $f'(\sigma) = f(\sigma)$
  અને $Rf(\sigma)w$ હોવાથી $Rf'(\sigma)f'(\sigma.n)$;
  તેથી $f'$ એ $P(\Gamma) \cup \{\sigma.n\}$નું અર્થઘટન છે.
  સ્પષ્ટ છે કે $\mSat/{M}{!B}[f'(\sigma.n)]$.
  બધા ઉપસર્ગ $\sigma' \in P(\Gamma)$ માટે
  $f(\sigma') = f'(\sigma')$ હોવાથી $\mSat{M}{\Gamma}[f']$.
  આમ $\mModel{M}, f'$ એ $\Gamma \cup
  \{\sFmla{\False}{!B}[\sigma.n]\}$ને સંતોષે છે.}
  }{}

\iftag{prvDiamond}{%
\item $\sFmla{\True}{\Diamond !B}[\sigma] \in \Gamma$ને
  $\TRule{\True}{\Diamond}$ લાગુ કરીને શાખા વિસ્તારીએ:
  \iftag{probDiamond}{અભ્યાસપ્રશ્ન.}{તેનાથી નવું !!{signed
      formula}~$\sFmla{\True}{!B}[\sigma.n]$ મળે છે, જ્યાં
    $\sigma.n$ શાખા પર નવો ઉપસર્ગ છે, એટલે
    $\sigma.n \notin P(\Gamma)$.
    \sourcecorrection{OLMD-056}{મૂળમાં આ નિયમનું નવું સૂત્ર $!A$
    સાથે લખાયું છે; પૂર્વપક્ષ $\Diamond !B$ અને આગળની સાબિતી
    મુજબ અહીં $!B$ હોવું જોઈએ.}
    $\Gamma$ સંતોષ્ય હોવાથી એવી $\Struct{M}$ અને
    $P(\Gamma)$નું અર્થઘટન~$f$ છે કે
    $\mSat{M}{\Gamma}[f]$, ખાસ કરીને
    $\mSat{M}{\Diamond !B}[f(\sigma)]$.
    $\Gamma \cup \{\sFmla{\True}{!B}[\sigma.n]\}$ સંતોષ્ય છે
    તે બતાવવાનું છે. તે માટે $P(\Gamma)
    \cup \{\sigma.n\}$નું અર્થઘટન નીચે મુજબ વ્યાખ્યાયિત કરીએ:

  $\mSat{M}{\Diamond !B}[f(\sigma)]$ હોવાથી એવું $w \in W$
  છે કે $Rf(\sigma)w$ અને $\mSat{M}{!B}[w]$. $f'$ એ $f$
  જેવું જ રાખો, માત્ર $f'(\sigma.n) = w$.
  $f'(\sigma) = f(\sigma)$ અને $Rf(\sigma)w$ હોવાથી
  $Rf'(\sigma)f'(\sigma.n)$; તેથી $f'$ એ
  $P(\Gamma) \cup \{\sigma.n\}$નું અર્થઘટન છે. સ્પષ્ટ છે કે
  $\mSat{M}{!B}[f'(\sigma.n)]$. બધા ઉપસર્ગ
  $\sigma' \in P(\Gamma)$ માટે $f(\sigma') = f'(\sigma')$
  હોવાથી $\mSat{M}{\Gamma}[f']$. આમ $\mModel{M}, f'$ એ
  $\Gamma \cup \{\sFmla{\True}{!B}[\sigma.n]\}$ને સંતોષે છે.}
\item $\sFmla{\False}{\Diamond !B}[\sigma] \in \Gamma$ને
  $\TRule{\False}{\Diamond}$ લાગુ કરીને શાખા વિસ્તારીએ:
  \iftag{probDiamond}{અભ્યાસપ્રશ્ન.}{તેનાથી કોઈ
    $\sigma.n \in P(\Gamma)$ માટે નવું !!{signed
      formula}~$\sFmla{\False}{!B}[\sigma.n]$ શાખા પર મળે છે
    (કારણ કે $\sigma.n$ વપરાયેલો હોવો જોઈએ). ધારો કે
    $\mSat{M}{\Gamma}[f]$, ખાસ કરીને
    $\mSat/{M}{\Diamond !B}[f(\sigma)]$. $f$ ઉપસર્ગોનું
    અર્થઘટન છે અને $\sigma$, $\sigma.n \in P(\Gamma)$ બંને
    હોવાથી $Rf(\sigma)f(\sigma.n)$ છે. તેથી
    $\mSat/{M}{!B}[f(\sigma.n)]$; એટલે $\mModel{M}, f$ એ
    $\sFmla{\False}{!B}[\sigma.n]$ને સંતોષે છે.}
  }{}
\end{enumerate}
\sourcecorrection{OLMD-057}{મૂળમાં બૉક્સ અને ડાયમંડના
નવા-ઉપસર્ગ કિસ્સાઓમાં $\Sat{M}{\Gamma}[f]$ વપરાયું છે;
અહીં વ્યાખ્યાયિત મોડલ-સંતોષ $\mSat{M}{\Gamma}[f]$ રાખ્યું છે.}
હવે બે શાખા આપતા સંભવિત અનુમાનો ધ્યાનમાં લઈએ.
\sourcecorrection{OLMD-058}{મૂળમાં આ નિયમોને ``બે પૂર્વપક્ષવાળા''
કહ્યા છે; દરેક નિયમને એક પૂર્વપક્ષ છે અને તે શાખાને બે ભાગમાં
વહેંચે છે, તેથી શાખાઓની સંખ્યા સ્પષ્ટ કરી છે.}
\begin{enumerate}
\item $\sFmla{\False}{!B \land !C}[\sigma] \in \Gamma$ને
  $\TRule{\False}{\land}$ લાગુ કરીને શાખા વિસ્તારીએ. તેનાથી
  બે શાખા મળે છે: ડાબી શાખા $\sFmla{\False}{!B}[\sigma]$માંથી
  અને જમણી શાખા $\sFmla{\False}{!C}[\sigma]$માંથી આગળ વધે છે.
  ધારો કે $\mSat{M}{\Gamma}[f]$, ખાસ કરીને
  $\mSat/{M}{!B \land !C}[f(\sigma)]$. ત્યારે
  $\mSat/{M}{!B}[f(\sigma)]$ અથવા
  $\mSat/{M}{!C}[f(\sigma)]$. પહેલા કિસ્સામાં
  $\mModel{M}, f$ એ $\sFmla{\False}{!B}[\sigma]$ને સંતોષે છે;
  એટલે ડાબી શાખા સંતોષ્ય છે. બીજા કિસ્સામાં
  $\mModel{M}, f$ એ $\sFmla{\False}{!C}[\sigma]$ને સંતોષે છે;
  એટલે જમણી શાખા સંતોષ્ય છે.
\item $\sFmla{\True}{!B \lor !C}[\sigma] \in \Gamma$ને
    $\TRule{\True}{\lor}$ લાગુ કરીને શાખા વિસ્તારીએ:
    અભ્યાસપ્રશ્ન.
\item $\sFmla{\True}{!B \lif !C}[\sigma] \in \Gamma$ને
    $\TRule{\True}{\lif}$ લાગુ કરીને શાખા વિસ્તારીએ:
    અભ્યાસપ્રશ્ન.
\end{enumerate}
\end{proof}

\begin{prob}
\olref[nml][tab][sou]{thm:tableau-soundness}ની સાબિતી પૂરી કરો.
\end{prob}

\begin{cor}
\ollabel{cor:entailment-soundness}
જો $\Gamma \Proves !A$ હોય, તો $\Gamma \Entails !A$.
\end{cor}

\begin{proof}
  જો $\Gamma \Proves !A$ હોય, તો કેટલાંક $!B_1$, \dots,
  $!B_n \in \Gamma$ માટે
  $\Delta = \{\sFmla{\False}{!A}[1], \sFmla{\True}{!B_1}[1],
  \dots, \sFmla{\True}{!B_n}[1]\}$નું બંધ !!{tableau} છે.
  બતાવવું છે કે $\Gamma \Entails !A$. એવું ન હોય તો કોઈ
  $\mModel{M}$ અને $w$ માટે બધા $i=1$, \dots,~$n$ પર
  $\mSat{M}{!B_i}[w]$ હોય, પરંતુ
  $\mSat/{M}{!A}[w]$. $f(1) = w$ રાખો. ત્યારે $f$ એ
  $P(\Delta)$નું~$\mModel{M}$માં અર્થઘટન છે અને
  $\mModel{M}$ એ~$f$ને અનુલક્ષીને~$\Delta$ને સંતોષે છે.
  પરંતુ \olref{thm:tableau-soundness} મુજબ બંધ !!{tableau}
  હોવાથી $\Delta$ અસંતોષ્ય છે—વિરોધાભાસ. તેથી
  $\Gamma \Entails !A$.
  \sourcecorrection{OLMD-059}{મૂળની છેલ્લી પંક્તિમાં ફરી
  $\Gamma \Proves !A$ નિષ્કર્ષ લખ્યો છે; આ ઉપપ્રમેય માટે જરૂરી
  નિષ્કર્ષ $\Gamma \Entails !A$ પુનઃસ્થાપિત કર્યો છે.}
\end{proof}

\begin{cor}
\ollabel{cor:weak-soundness}
જો $\Proves !A$ હોય, તો $!A$ બધાં નિદર્શનોમાં સત્ય છે.
\end{cor}

\end{document}
